%-------------------------
% Resume for Kiron Sukesan
% Version: PhD / faculty outreach (AI alignment). Keep a separate industry version with the recruiter-facing summary.
%-------------------------
\documentclass[10pt, letterpaper]{article}

\usepackage[T1]{fontenc}
\usepackage[top=0.5in, bottom=0.5in, left=0.65in, right=0.65in]{geometry}
\usepackage{enumitem}
\usepackage{titlesec}
\usepackage{hyperref}
\usepackage{xcolor}
\usepackage{tabularx}
\usepackage{array}

% Colors
\definecolor{accent}{RGB}{30, 64, 120}
\definecolor{lightgray}{RGB}{130, 130, 130}

% Hyperlink setup
\hypersetup{
  colorlinks=true,
  urlcolor=accent,
  linkcolor=accent
}

% Section formatting
\titleformat{\section}
  {\bfseries\normalsize\color{accent}}
  {}{0em}{}
  [\vspace{-4pt}\textcolor{accent}{\rule{\linewidth}{0.8pt}}\vspace{2pt}]

\titlespacing*{\section}{0pt}{7pt}{4pt}

% No indentation
\setlength{\parindent}{0pt}

% Compact list style
\setlist[itemize]{
  leftmargin=1.2em,
  itemsep=1.5pt,
  parsep=0pt,
  topsep=2pt,
  label=\textbullet
}

% Custom commands
\newcommand{\resumeheader}[5]{%
  \begin{center}
    {\LARGE\bfseries #1}\\[4pt]
    {\color{lightgray}\small
      #2 \,$\cdot$\, #3 \,$\cdot$\, \href{#4}{linkedin.com/in/kironsukesan} \,$\cdot$\, \href{#5}{github.com/kironsuk}
    }
  \end{center}
}

\newcommand{\experienceentry}[4]{%
  \vspace{4pt}
  \noindent
  \begin{tabularx}{\linewidth}{@{}X r@{}}
    \textbf{#1} \textit{\color{lightgray}at} \textbf{#2} & \textcolor{lightgray}{\small #3} \\
    \multicolumn{2}{@{}l@{}}{\textit{\color{lightgray}\small #4}}
  \end{tabularx}
  \vspace{1pt}
}

\newcommand{\educationentry}[4]{%
  \vspace{2pt}
  \noindent
  \begin{tabularx}{\linewidth}{@{}X r@{}}
    \textbf{#1} & \textcolor{lightgray}{\small #3} \\
    \textit{#2} & \textcolor{lightgray}{\small #4}
  \end{tabularx}
}

\begin{document}
\pagestyle{empty}

%------------------------
% HEADER
%------------------------
\resumeheader
  {Kiron Sukesan}
  {sukesan@umich.edu}
  {Ann Arbor, MI}
  {https://linkedin.com/in/kironsukesan}
  {https://github.com/kironsuk}

\vspace{5pt}

%------------------------
% SUMMARY
%------------------------
\section{Summary}
Research assistant in Wei Lu's group at the University of Michigan (EECS), testing a state-space model of hippocampal time cells against neural recordings. Previously nine years in industry: production ML and computer vision at Compass, then distributed storage at AWS, including two years managing the DynamoDB Storage Node team. Research background in Bayesian active learning and closed-loop BCIs, with three peer-reviewed publications. Applying to PhD programs for Fall 2027 to work on understanding what large models represent internally and how multiple AI systems behave together.

%------------------------
% SKILLS
%------------------------
\section{Skills}
\begin{tabularx}{\linewidth}{@{} l X @{}}
  \textbf{Languages}      & Python, Java, C++, Rust, JavaScript, MATLAB, Bash \\[2pt]
  \textbf{ML / AI}        & PyTorch, Deep Learning, Reinforcement Learning, Computer Vision, Gaussian Processes, Bayesian Inference, Semi-Supervised Learning, ANN/HNSW \\[2pt]
  \textbf{Storage Systems}        & Distributed Database Internals, Consensus Protocols (Paxos, Raft), Database Engine Implementation (LSM tree and B-tree)  \\[2pt]
\end{tabularx}

%------------------------
% EXPERIENCE
%------------------------
\section{Experience}

\experienceentry{Research Assistant}{Wei Lu Group, EECS (ECE), University of Michigan}{Oct 2026 -- Present}{Ann Arbor, MI}
\begin{itemize}
  \item Testing a state-space model of hippocampal time cells (Lu et al., arXiv 2507.13638) against calcium-imaging recordings of mouse CA1, asking whether the model's time-cell-like units match recorded neurons
  \item Building the model-to-data comparison pipeline in Python: matched task environment, the same time-cell classifier applied to model units and recorded neurons, and a similarity metric compared across initializations (HiPPO, random, real-diagonal)
\end{itemize}

\experienceentry{Lead Engineer}{Freelance}{Apr 2026 -- Present}{Ann Arbor, MI}
\begin{itemize}
  \item Architected and developed AppMan, a full-stack apartment management platform using Angular, PostgreSQL, and Docker, modernizing legacy Oracle APEX workflows for property operations, payments, leasing, taxes, and task management. In production at a 1,000+-unit property management company.
  \item Designed and built a test-driven engineering infrastructure with unit, integration, adversarial, migration, authorization, and smoke tests.
\end{itemize}


\experienceentry{Software Development Manager}{Amazon Web Services (AWS)}{Aug 2024 -- Apr 2026}{DynamoDB Storage Node, Seattle, WA}
\begin{itemize}
  \item One of the engineering leaders for DynamoDB's Storage Node, among AWS's highest-scale distributed storage systems, overseeing architecture, delivery, and team health. Specifically owned the TTL system, replication protocol, general correctness of DynamoDB Storage, and feature expansion efforts within storage.
  \item Drove hiring, mentoring, and performance management for a team of 10+ junior, mid-level and senior software engineers; established engineering standards and oncall culture.
\end{itemize}

\experienceentry{Senior Software Engineer}{Amazon Web Services (AWS)}{Jan 2023 -- Aug 2024}{DynamoDB Storage Node API Team, Seattle, WA}
\begin{itemize}
  \item Designed and implemented features for the DynamoDB Storage Node API layer, contributing to a system serving millions of customers with stringent latency, availability, and durability requirements.
  \item Led oncall and operational excellence initiatives, driving root cause analyses and systemic improvements to reduce customer-impacting events across the storage tier.
\end{itemize}

\experienceentry{Machine Learning Engineer II / Senior Machine Learning Engineer}{Compass}{Dec 2019 -- Nov 2022}{Seattle, WA}
\begin{itemize}
  \item Built and owned computer-vision-based ML models used by real estate agents across the platform, directly influencing transaction outcomes.
  \item Architected end-to-end ML pipelines spanning data ingestion, feature engineering, model training, evaluation, and production deployment on cloud infrastructure.
\end{itemize}

\experienceentry{Software Development Engineer I/II}{Amazon Kindle}{Oct 2017 -- Nov 2019}{Greater Seattle Area}
\begin{itemize}
  \item Designed and shipped features for Amazon Kindle's internal data processing pipeline, collaborating with senior engineers on distributed systems design, API contracts, and service reliability.
  \item Promoted to SDE II after demonstrating technical ability with an end-to-end solution for Kindle compliance with GDPR's requirement to return data to customers on request.
\end{itemize}

\experienceentry{Data Scientist / Research Engineer}{Washington University in St.\ Louis}{May 2016 -- Aug 2017}{Laboratory of Sensory Neuroscience and Neuroengineering, St. Louis, MO}
\begin{itemize}
  \item Developed a Bayesian active learning framework using Gaussian Process Classification to efficiently estimate audiometric threshold surfaces, reducing clinical measurement burden while preserving diagnostic accuracy.
  \item Designed a novel ML-based hardware calibration algorithm for audiometric equipment, improving signal fidelity and measurement consistency across experimental sessions.
  \item Co-authored three peer-reviewed publications on Bayesian active learning, psychometric field estimation, and online machine-learning audiometry.
\end{itemize}

\experienceentry{Software Developer}{Purnatva Solutions Pvt.\ Ltd.}{Sep 2014 -- Aug 2015}{Bengaluru, India}
\begin{itemize}
  \item Implemented real-time GPS bus tracking module for the Falcon Mobile App, supporting large-scale deployment with 10,000+ active users across the platform.
\end{itemize}

\experienceentry{Biomedical Engineer, Undergraduate Research}{Wash.\ U.\ in St.\ Louis}{Aug 2013 -- May 2014}{Department of Biomedical Engineering}
\begin{itemize}
  \item Designed and built a wearable EEG recording device for overnight clinical use, using custom analog circuitry, embedded software, and mechanical design in close collaboration with two biomedical engineering peers and clinical partners.
  \item Validated device performance against clinical EEG standards and iterated on hardware and firmware based on clinician feedback.
\end{itemize}

\experienceentry{Neurosurgery Research Intern}{University of Washington}{Jun 2013 -- Aug 2013}{Brain-Computer Interface Lab (Dr.\ Jeff Ojemann), Seattle, WA}
\begin{itemize}
  \item Developed a signal analysis algorithm for electrocorticographic (ECoG) neural signals; diagnosed and resolved electromagnetic noise artifacts that had been corrupting recorded data.
  \item Constructed a closed-loop BCI device that enabled motor-impaired patient hand movement through real-time classification of theta wave activity from ECoG recordings.
\end{itemize}

%------------------------
% EDUCATION
%------------------------
\vspace{2pt}
\section{Education}

\educationentry{Washington University in St.\ Louis}{Master of Science, Computer Science}{2015 -- 2016}{St. Louis, MO}
\vspace{6pt}
\educationentry{Washington University in St.\ Louis}{Bachelor of Science, Biomedical Engineering}{2010 -- 2014}{St. Louis, MO}

%------------------------
% PUBLICATIONS
%------------------------
\section{Publications}
\begin{itemize}
  \item Song, X. D., \textbf{Sukesan, K. A.}, \& Barbour, D. L. ``Bayesian active probabilistic classification for psychometric field estimation.'' \textit{Attention, Perception, \& Psychophysics}, 80 (2018), 798--812. DOI: \href{https://doi.org/10.3758/s13414-017-1460-0}{10.3758/s13414-017-1460-0}.
  \item Barbour, D. L., DiLorenzo, J. C., \textbf{Sukesan, K. A.}, Song, X. D., Chen, J. Y., Degen, E. A., Heisey, K. L., \& Garnett, R. ``Conjoint psychometric field estimation for bilateral audiometry.'' \textit{Behavior Research Methods}, 51 (2019), 1271--1285. DOI: \href{https://doi.org/10.3758/s13428-018-1062-3}{10.3758/s13428-018-1062-3}.
  \item Barbour, D. L., Howard, R. T., Song, X. D., Metzger, N., \textbf{Sukesan, K. A.}, et al. ``Online machine learning audiometry.'' \textit{Ear and Hearing}, 40(4) (2019), 918--926. DOI: \href{https://doi.org/10.1097/AUD.0000000000000669}{10.1097/AUD.0000000000000669}.
\end{itemize}

\end{document}
